\section{Introduction}
\label{sec:rl_motivation}
\rev{
In the preceding chapters, we established the conceptual and formal basis for this dissertation. Novelty arises when an encounter exceeds an agent's existing knowledge or competence, and the structure of the agent's abstractions helps determine whether the resulting mismatch can be localized and addressed. In \cref{ch:benchmarks}, we then presented environments and metrics for studying how such mismatch affects task performance. This chapter turns to the accommodation side of the dissertation's central argument: \emph{how can an agent expand its behavioral repertoire when novelty makes an action required by its task no longer work?}

For a goal-directed agent, incomplete knowledge can become apparent in two different ways. The agent may be unable to construct a plan, indicating that its symbolic model lacks the structure needed to reach the goal. Alternatively, it may construct a valid plan but fail while executing one of its actions. In the second case, the plan can still describe what must be achieved even though the agent no longer knows how to achieve it. The neurosymbolic architecture presented in this chapter addresses both forms of failure, while the chapter's primary learning problem concerns the second: recovery from obstructive novelty that invalidates an existing behavior.

The action abstractions introduced in \cref{sec:action_abstractions} provide the structure needed for this recovery. Each symbolic operator specifies the conditions under which an action applies and the effects it is intended to produce; an associated executor specifies how the agent realizes that action in the environment. If the operator's preconditions hold but its expected effects do not occur, the mismatch can be localized to the active operator--executor pair. The agent can first update its knowledge, replan, or test inexpensive symbolic hypotheses. If these strategies are insufficient, it can learn a new executor for the affected operator rather than relearning the entire task. Operators and executors unaffected by the novelty remain available for reuse.

In the terminology developed in \cref{ch:novelty_open_world}, this is \emph{localized behavioral accommodation}. The agent changes the part of its behavioral structure that no longer matches the environment, expanding its repertoire at the point of failure while preserving competence that remains valid. The claim is deliberately scoped: the agent does not construct an arbitrary ontology from scratch, but adapts within a symbolic task interface whose operators, state descriptions, and executor relationships provide the structure for diagnosis and recovery.

We instantiate that idea at two levels. First, we present a neurosymbolic cognitive architecture for open-world novelty handling~\citep{goel2024neurosymbolic}. The architecture combines symbolic and neural inference with a cascading recovery process that connects novelty detection and characterization to an appropriate response. Lightweight symbolic strategies are attempted first; learning is invoked when the existing model and behaviors are insufficient. Second, we present \emph{\RAPidLearn{}} (Learning to Recover and Plan Again)~\citep{goel2022rapidlearn}, the architecture's executor-learning mechanism. \RAPidLearn{} formalizes executor discovery as a targeted reinforcement learning problem defined by the failed operator, its intended effects, and the states from which planning can resume. Existing knowledge and behaviors guide exploration, and a successful executor is retained for future use.

The axe novelty from \cref{sec:benchmarks_motivation} illustrates the process. The agent has a plan that requires it to break a tree, but after novelty is introduced, the usual \texttt{break} executor no longer produces a tree log unless the agent uses an axe. The failed effect identifies both where recovery is needed and what recovery must accomplish. Rather than learning the crafting task again, the agent learns a replacement executor that uses the axe---obtaining it first when necessary---to restore the intended result. It can then resume planning with the remainder of its knowledge and behavior unchanged.

In the remainder of the chapter, we develop and evaluate these two contributions. \Cref{sec:architecture} presents the neurosymbolic architecture and its detection, characterization, and recovery mechanisms. \Cref{sec:ipt} formalizes the integrated planning task and executor discovery problem, and \cref{sec:rapid_learn_method} presents the \RAPidLearn{} algorithm and its knowledge-guided exploration strategies. Finally, in \cref{sec:rl_evaluation}, we evaluate \RAPidLearn{} on five novelty scenarios in \NovelGridworlds{} and the integrated architecture across 288 novelties spanning ten categories in Polycraft.}
